\section{Related Work}
\label{sec:related}
\paragraph{Secure voice and media-path authentication.}
ZRTP establishes key agreement in the media path for secure RTP, including
short authentication strings and key continuity~\cite{rfc6189}. Krasnowski,
Lebrun, and Martin study authenticated key exchange over voice-only channels,
using signatures and short authentication strings with symbolic
verification~\cite{voicekeys}. Their work is especially close evidence that
cryptographic exchange through an audio domain is not itself new. Our different
target is a participating software agent and business executor using registered
keys and scoped customer authorization. We have not demonstrated the codec
robustness or formal guarantees of that prior work, and human short-string
verification is not our authentication mechanism.

Classical modems establish the feasibility of transporting bits through
telephone infrastructure; V.90 also illustrates how achievable rates depend
on the digital/analog path~\cite{v90}. \sys{}'s exact-level PCM profile is a
controlled digital-path experiment, not an implementation of V.90 or evidence
that arbitrary modern speech codecs preserve such signaling. Its FSK profile
is deliberately simple. A serious networking contribution needs to measure
these adaptation constraints and compare appropriate coding and recovery
baselines, rather than claim that wrapping a familiar transport is sufficient.

\paragraph{Protected transports and delegation.}
TLS 1.3 and QUIC supply established cryptographic transport and recovery
machinery~\cite{rfc8446,rfc9000,rfc9001,rfc9002}. We reuse them rather than
invent a cipher or key exchange. Macaroons show how context and attenuated
delegation can be encoded in authorization credentials~\cite{macaroons}.
OAuth certificate-bound tokens and DPoP demonstrate standardized possession
binding~\cite{rfc8705,rfc9449}. \sys{} combines analogous separation of
identity and permission with a concrete audio-carried endpoint. It does not
claim its custom grants improve those standards cryptographically.

PAP, A2A, and MCP provide complementary agent/application
protocols~\cite{pap03,a2a,mcp}. Their scopes and evolving drafts matter when
making an interoperability claim. The repository's PAP adapter is a local,
custom extension of its earlier supported profile, and the current public
PAP specification should not be characterized as omitting authentication,
authorization, or prompt-injection considerations.

\paragraph{Perceptual and prompt injection.}
DolphinAttack exploits microphone nonlinearity to introduce inaudible voice
commands, and Light Commands demonstrates optical injection into
microphones~\cite{dolphin,lightcommands}. These are attacks on perceptual input,
not cryptographic packet forgery. Indirect prompt injection exposes the risk
of interpreting external data as instructions in tool-using
applications~\cite{greshake}. AudioHijack studies auditory prompt injection
against audio-language models~\cite{audiohijack}. \sys{} places an independent
authorization gate on a subset of downstream effects. It does not stop
perceptual injection or establish a general boundary between all data and
instructions inside a model. The relevant comparison is which effects remain
possible after manipulation, under explicitly matched tool permissions.
